%!TEX program = uplatex
\RequirePackage{plautopatch}

\documentclass[uplatex,dvipdfmx]{jlreq}

% 余白調整（1ページ収め）
\usepackage[top=18mm,bottom=20mm,left=20mm,right=20mm]{geometry}

\usepackage{amsmath,amssymb}
\usepackage{url}

\pagestyle{empty}

% 行間少し詰める
\renewcommand{\baselinestretch}{0.96}

%--------------------------------
% 講演マクロ
%--------------------------------

\newcommand{\talk}[3]{%
\noindent
#1\quad #2\par
\hspace*{2em}#3\par\smallskip
}

%--------------------------------
% タイトル情報
%--------------------------------

\newcommand{\EwMnumber}{第3回}
\newcommand{\EwMtitle}{粘性解理論への招待\\[2mm]\large 幾何系・代数系の人たちにも知ってもらいたい}
\newcommand{\EwMdate}{1997年5月16日（金）14:30 ～ 5月17日（土）17:00}
\newcommand{\EwMplace}{東京都文京区春日1-13-27 中央大学理工学部5号館 5533号室}

%--------------------------------
% タイトル表示
%--------------------------------

\newcommand{\EwMtitleblock}{

\vspace*{-5mm}

\begin{center}

{\Large ENCOUNTER with MATHEMATICS}

\vspace{2mm}

{\large 数学との遭遇}

\vspace{4mm}

{\Large \bfseries \EwMnumber\ \EwMtitle}

\vspace{4mm}

\EwMdate

\vspace{2mm}

於：\EwMplace

\end{center}

\vspace{3mm}


\vspace{2mm}

}

\begin{document}

\EwMtitleblock

%--------------------------------
% プログラム
%--------------------------------

\section*{プログラム}

\subsection*{1日目（5月16日（金））}

\talk
{14:30～15:40}
{粘性解への招待 I}
{石井 仁司 氏（都立大・理）}

\talk
{16:20～17:30}
{粘性解という弱解の幾何学的側面 I}
{儀我 美一 氏（北大・理）}

\subsection*{2日目（5月17日（土））}

\talk
{10:30～11:30}
{粘性解への招待 II}
{石井 仁司 氏（都立大・理）}

\talk
{11:50～12:50}
{粘性解という弱解の幾何学的側面 II}
{儀我 美一 氏（北大・理）}

\talk
{14:20～15:20}
{「半連続粘性解」について}
{小池 茂昭 氏（埼玉大・理）}

\talk
{15:50～16:50}
{粘性解と制御理論の関わり}
{長井 英生 氏（阪大・基礎工）}

%--------------------------------
% 案内文
%--------------------------------

\bigskip

別掲の趣旨に沿った集会の第3回を以上のような予定で開催いたします。  
非専門家向けに入門的な講演をお願い致しました。  
代数系・幾何系の方々の多くのご参加をお待ちしております。  
講演者による講演内容へのご案内を別紙に添付いたしますのでご覧下さい。

%--------------------------------
% 連絡先
%--------------------------------

\section*{連絡先}

112 東京都文京区春日1-13-27 中央大学理工学部数学教室

tel : 03-3817-1745

ENCOUNTER with MATHEMATICS  
e-mail : encmath@math.chuo-u.ac.jp  

三松 佳彦 : e-mail : yoshi@math.chuo-u.ac.jp  

\url{http://www.math.chuo-u.ac.jp/ENCwMATH}

\end{document}
